For a trigonometric polynomial $f(\theta)=\sum_{k=0}^n c_ke^{ik\theta}$ whose zeros are real, normalized by $\max_{\theta\in[0,2\pi]}|f(\theta)|=1$, the requested estimate is
\[
 \int_0^{2\pi}|f(\theta)|\,d\theta\le4.\tag{1}
\]
The current problem page records an affirmative solution by Kristiansen and, for complex coefficients, by Saff and Sheil-Small. This attempt proves sharp special families, not their general theorem.

There is a degenerate case to exclude when reading the displayed statement literally: $f(\theta)=e^{im\theta}$ has no zeros, so ``all zeros are real'' is vacuous, but its integral is $2\pi$. Accordingly, the analysis below concerns the nondegenerate zero-bearing case. After removing a nonvanishing factor $e^{im\theta}$, this is a polynomial $Q(e^{i\theta})$ with $Q(0)\ne0$, $\deg Q=d\ge1$, all of whose algebraic zeros lie on the unit circle. A nonzero algebraic zero $\rho e^{it}$ corresponds to trigonometric zeros $\theta=t-i\log\rho+2\pi j$, which explains this correspondence. The degenerate observation is a statement convention, not a claimed new disproof of the classical result.
